\chapter{群論}
\section{基礎事項}
\subsection{位数}

\begin{prop} $G$を有限群とする. 
\ben 
\item $|G|$が素数$p$で割り切れるとき, 位数が$p$の元が存在する. 
\een 
\end{prop}

\begin{prac}
$\sigma = \bep 1&2&3&4&5&6&7&8&9\\ 2&1&4&5&3&7&8&9&6 \enp\in S_9$の位数を求めよ.
\end{prac}

\tbox{ 
\begin{egprob}
$G$は位数$n$の有限群で, 次の性質を満たす. 
\begin{center} 
$n$の任意の約数$d>0$に対し, $G$は位数$d$の部分群をただ一つ持つ. 
\end{center}
$G$はどのような群か. 
\end{egprob}
}{title=分からなかった問題(H10京大数学I)}
    \begin{egans} \href{https://math.stackexchange.com/questions/1302635/a-finite-group-which-has-a-unique-subgroup-of-order-d-for-each-d-mid-n}{参考}
位数$d$の元があるものとする. $x,y$を位数$d$の元とすると, 条件より$\lan x\ran = \lan y\ran$なので, $\lan x\ran$には位数$d$の元がすべて含まれる. $\lan x \ran \cong \mb{Z}/d\mb{Z}$なので$\lan x\ran$の中には位数$d$の元がちょうど$\phi(d)$個存在する. よって$G$には位数$d$の元が$\phi(d)$個存在するか, 1個も存在しない. よって
\[|G|=n \leq \sum_{0<d|n} \phi(d) \]
となる. \tf{ここで, 右辺は必ず$n$に等しいため}すべての$d$に対して位数$d$の元は$\phi(d)$個持たなければならない. よって示せた. \qed 
\end{egans}
テクニカルだが, 次のような使い方も覚えておきたい. 
\begin{egprob}
$G$を位数$30$の群とする. Sylow$p$部分群の個数を$n_p$とするとき, $n_3, n_5$のどちらかは1であることを示せ. 
\end{egprob}
\begin{egans}
Sylowの定理の主張を見て貰うと, $n_3 \in \set{1,10}$, $n_5\in \set{1,6}$である. よって$n_3=10, n_5=6$として矛盾を導けばよい. もしこうなっていたら, $G$には位数3の元が20個はあり, 位数5の元が24個あることになる. よっておかしい. \qed 
\end{egans}


\begin{egprob} (H28基礎科目II,4) 
$A= \sumib{0&-1\\ 1&1}$, $B=\sumib{0&1\\1&0}$で生成される$\mr{GL}_{2}(\mb{C})$の部分群$G$について, 以下の問に答えよ. 
\ben 
\item $G$の位数を求めよ. 
\een 
\end{egprob}




\subsection{有限Abel群}
\begin{egprob} (H26基礎科目II, 4)
$G = \cyc{4}\times \cyc{6}\times \cyc{9}$の指数3の部分群の個数を求めよ. 
\end{egprob}
\begin{egans}
指数3の部分群$H$は$3G = \cyc{4}\times 3\mb{Z}/6\mb{Z}\times 3\mb{Z}/9\mb{Z}$を含む. $G/3G = \cyc{3}\times \cyc{3}$だから, $H$は$G/3G$の位数3の部分群に対応する. これは位数3の元の個数の半分なので, $4$個である. 
\end{egans}

\begin{egprob} (H25数学I, 3) \label{prob:H25-I-3}
$p$は素数とする. アーベル群$A$は位数$p^4$であり, 位数$p$の部分群$N$で$A/N\cong \mb{Z}/p^3\mb{Z}$となるものを持つとする. このような$A$を同型を除いてすべて求めよ. 
\end{egprob}

\begin{egans}
仮定より$A$には位数$p^3$以上の元が存在する. もし位数$p^4$の元が存在すれば$A\cong \cyc{p^4}$である. 存在しないとする. 有限アーベル群の構造定理から$A\cong \cyc{p} \times \cyc{p^3}$となるしかないが, 実際この群は$N = \cyc{p}\times \set{0}$という部分群によって$A/N \cong \cyc{p^3}$となるのでよい. よって
\[A\cong \cyc{p^4}, \quad \cyc{p}\times \cyc{p^3}. \]
\end{egans}






\newpage 
\subsection{部分群・正規部分群}
\tbox{
\begin{prop}
\ben 
\item $\phi:G\to H$が群準同型なら$\ker{\phi}$は正規部分群である. 
\item $G$の中心$Z(G)$は正規部分群である. 
\item \tf{正規部分群は共役類の合併である. }
\item 指数$2$の部分群は必ず正規である. とくに$A_n\norgr S_n$である. (より一般に, $G$の位数を割り切る最小の素数$p$に対し, 指数$p$の部分群は正規である\footnote{こちらも証明は知っておいた方が良い. しかし指数2のほうは答案にすぐ使ってもよいくらい「常識」であろう. }. )
\item $G$を群, $H,K$は$G$の部分群, $N$は$G$の正規部分群とするとき, $HK$が部分群であることと$HK=KH$が成り立つことは同値. とくに, 正規部分群との積 $HN$は必ず$G$の部分群である. 
\item (\tf{重要}) $N_1,N_2$を$G$の正規部分群で$N_1\cap N_2 = \set{1_G}$かつ$G=N_1N_2$のとき, $G\cong N_1\times N_2$である. 
\een
\end{prop}
}{title=正規部分群のまとめ}
\prf (1),(2):Very Easy. \\
(3): $H\norgr G$とする. $ghg^{-1}\in H$なので$h\in H$の共役元はすべて$H$に属す. そこから明らかであろう. \qed \\
(4): $H$が指数2ならすべての$g$に対して$gH\in \set{H,Hg} $である. もしある$g$で$gH \neq Hg$なら$gH = H$であり$g\in H$だが$gH = H = Hg$でおかしい. \qed \\
(5): $hkh'k'$を$HK$の元と見たければ$HK=KH$を使うであろうことから前半が分かる. 後半も$HN=NH$は容易に分かるので前半の主張から言える. \qed \\
(6): (5)より$N_1N_2$は部分群. $\phi:N_1\times N_2 \to G$を$\phi(n_1,n_2) = n_1n_2$で定めて,全単射準同型であることを示せばよい. 単射であることはどう示せるか考えて見よ. \qed 
\begin{prac}
(6)の証明を埋めよ. 
\end{prac}








\begin{eg}
$G$が位数$pq$ ($p,q$は異なる素数)で, Sylow部分群$S_p,S_q$が正規であることが分かっているとする. このとき
$G = S_pS_q$と$S_p\cap S_q=\set{1_G}$が満たされるので$G\cong S_p\times S_q \cong C_{pq}$ ($C_n = \mb{Z}/n\mb{Z}$)である. 実は, $p<q$で$p$が$q-1$を割らないなら, $G\cong C_{pq}$であることが言える(これもそれなりに有名である). (See ..)
\end{eg}

正規部分群に関連して, 正規化群というものがある. 
\begin{defn}
群$G$とその部分群$H$に対し, $H$の\tf{正規化群}$N_{G}(H)$とは, 
\[N_{G}(H) = \set{\eta\in G \mid \eta H \eta^{-1} = H}\]
により定義される群である. すなわち, 共役コセット集合$\set{gHg^{-1} \mid g\in G}$に$G$が共役で作用するときの$H$の固定部分群である. 
\end{defn}

\begin{tips} H29専門科目3番の体論のように, 正規化群をわざわざ書いてくれることもある.  とはいえやっぱり正規化群くらいは覚えておいてもらいたい. 
\end{tips}

\begin{egprob}
$S^n$の指数2の部分群は$A_n$に限るか. 
\end{egprob}

\begin{egans} 指数2の部分群$N$とする. $N$は正規である. $G/N\cong \mb{Z}/2\mb{Z}$なので全射準同型$S_n \to \mb{Z}/2\mb{Z}$がある. 逆に全射準同型$\phi : S_n \to \mb{Z}/2\mb{Z}$を与えると準同型定理より$\ker{\phi}$は指数2の正規部分群である. よって, $N$は$\ker{\phi}$として必ず書ける. さて, $S_n$は互換で生成され, 互換同士はサイクルの型が同じであるから共役である. もしある互換$\sigma$が$\phi(\sigma) = 0$を満たせば, 任意の互換$\sigma' = \alpha \sigma \alpha^{-1} $に対し
\[\phi(\sigma') = \phi(\alpha) + \phi(\sigma) + \phi(\alpha^{-1}) = \phi(\sigma) = 0\quad \mr{in} \mb{Z}/2\mb{Z}\]
と分かる. 逆に$\phi(\sigma) = 1$なら$\phi(\sigma') = 1$. よって$\phi$は自明な準同型か符号関数しか存在しない. そのうち全射であるものは符号関数で, その核は$A_n$である. \qed 
\end{egans}

部分群からそれなりに大きい正規部分群を構成したいときに次が役立ったことがある. 
\begin{egprob}
$G$を群, $H$を部分群とする. $H$に含まれる$G$の最大の正規部分群は$K=\bigcap_{g\in G}g^{-1}Hg$であるか? 
\end{egprob}
\begin{egans}
$K$が正規部分群であることと$N$が$H$に含まれる$G$の正規部分群なら$N\subset K$を示せばよい.\par 
まず$K$は部分群. $x\in K$とすると$y^{-1}xy$は$(gy)^{-1}H(gy)$に含まれ, $g$がすべて動くので$g^{-1}Hg^{-1}$に含まれる. よって$y^{-1}xy\in K$であり正規である. また, $N\subset K$は正規部分群の定義から明らかである. 
\end{egans}
\begin{egans}
$K$はcosetへの作用$G\to \mr{Aut}(G/H)$の核なので. 
\end{egans}


\begin{egprob}
部分群をちょうど$5$個持つような有限群$G$の構造を決定せよ. ただし, $G,\set{1}$も部分群として数える. 
\end{egprob}

\begin{egans} (H3数学II-1)
巡回群$C_n$から考えると, 部分群と$n$の約数は対応する. よって$n=p^4$という形である. 
シロー部分群を考えれば, $|G|$は4つ以上の素因数を持たない. 
$|G|$が素数$p<q<r$で割れるとき, $G$はアーベル群でなければならない. なぜなら, $S_p,S_q,S_r$がどれも一つだけであるから. $S_{p},S_q,S_r$はもはや自明な部分群を含んではならないので, $S_p,S_q,S_r$は巡回群である. よって$G\cong C_p\times C_q\times C_r$と書けるが, これは実際には条件を満たさない.\\
$|G|$の素因数が$p<q$のみであるとする. $G$が巡回群でないアーベル群の場合, $C_{p^i}\times C_{p^jq^k}$という形か, これの$p,q$を入れ替えた形のものが考えられる. ただし$i,j,k\geq 1$である. このとき, $C_{p^jq^k}$に部分群が4個はあるので, この形では部分群が8個は見つかり, 不適. 以降非アーベル群から考える. Sylow$p,q$部分群の個数をそれぞれ$n_p$,$n_q$とする. $n_p + n_q \leq 3$が必要である. $p<q$なので$n_q = 1$である. 一方, $n_p >1$であれば$n_p \geq q-1\geq 2$であるから, $q\geq 5$のときは$n_q=1$となり, $G$がアーベル群になるから不適. よって$q=3$で, 非アーベル群なので$G=S_3$が考えられる. これは互換で生成される位数2の部分群を3つ, 位数3の部分群$A_3$を持つので, 不適である. \\
$G$が$p$群かつ巡回群でないとする. $G=C_p\times C_p$のとき, 部分群が5個であるのは$p=2$のときである. $|G|\geq p^3$とする. 位数$p$の元が$c$個あるとすれば, 位数$p$の部分群がちょうど$c/(p-1)$個ある. $c=p^3-1$なら$c/(p-1) = p^2+p+1 > 3$なので不適. よって位数$p^2$の元$g$が存在する. $\lan g \ran $に属さない$h$を取ると, $\lan h \ran, \lan g \ran, \lan g^p \ran $は相異なる部分群である. これらの合併は$p^2 + p$個の元からなるが, $|G|\geq p^3 > p^2 + p$だから, この合併に属さない元$s$があり, $\lan s\ran$はまた相異なる部分群なので, 不適である. 以上より
\[C_{p^4}, C_{2}\times C_{2}, \]
\end{egans}



\newpage 
\subsection{共役類}

\tbox{ 
\begin{prop}
\ben 
\item $\maf{S}_{d}$の共役類と\tf{自然数$d$の分割}は対応する. 
\een
\end{prop}
}{title=共役類まとめ}
\begin{eg}
$\maf{S}_{3}$の共役類は$\set{1},\set{(12),(23),(13)},\set{(123),(132)}$である.  $\maf{S}_{4}$の共役類は$\set{1},\set{互換}$,$\set{\text{3-cycle}}$, $\set{\text{4-cycle}}$, $\set{\text{位数2偶置換}}$である. 個数はそれぞれ1,6,8,6,3である. また, $\maf{S}_{4}$のこの情報から「指数$3$の正規部分群が存在しない」ことが導ける. [証明]指数$3$なら位数$8$である. 正規部分群は共役類の合併であるから, 共役類をつなぎ合わせて$8$個の元を持つようにできるかを考えたいが, そのような方法は一つもない. よって存在しない. \footnote{実は$\maf{S}_{5}$の自明でない正規部分群は$A_5$のみであることが知られている. $\maf{S}_{d}$ ($d\neq 4$)でも同様. $d=4$の場合, Kleinの四元群というものがある. }. 
\end{eg}
\begin{egprob}
\end{egprob}



\begin{prac}
$S_5$の共役類とそれぞれの元の個数をすべて答えよ. (7個ある)
\end{prac}

\begin{theo}
$G$を有限群とする。$G$の共役類の完全代表系で中心$Z(G)$の元でないものを$\{ x_1,\cdots, x_n\}$とするとき, 
\[|G| = |Z(G)| + \disp\sum_{i} |C(x_i)| = |Z(G)|  +  \disp\sum_{i} (G:C_{G}(x_i))\]
である。ただし, $C_{G}(x_i)$は中心化群$\{ g\in G\mid gx_i = x_i g\}$である。
\end{theo}
$Z(G)$および$(G:C_{G}(x_i))$はLagrangeの定理より$|G|$の約数であるから, 類等式の数論的制約によってその構造が決定されることがある。
\begin{eg}
$p$を素数とするとき, 位数$p^2$の群$G$はAbel群, すなわち$|Z(G)| = p^2$であることを示す。類等式から$|Z(G)| = 1$ではありえない。$|Z(G)| = p$であるとして矛盾を導く。このとき$Z(G)$と$G/Z(G)$はともに位数$p$の巡回群である。生成元$a\in Z(G)$と$[b]\in G/Z(G)$をひとつ選ぶ(したがって$a\neq 1, b\notin Z(G)$)。$g\in G$を任意に取る。$[g] = [b]^{n}$となる自然数$n$が存在する。このとき$b^{n} g^{-1}\in Z(G)$だから$b^{n} g^{-1} = a^{m}$となる自然数$m$が存在する。よって$g=a^{-m} b^{n}$だから$G=\lan a,b \ran$である。すなわち, $G$の任意の元は$zb^{n}$と書ける。$z_1,z_2\in Z(G)$, $n_1,n_2\in \mb{N}$に対し, $z_{1}b^{n_1}$と$z_{2}b^{n_2}$が可換であることは容易に分かるので, $G$のすべての元は可換になり矛盾である。よって$G=Z(G)$である。
\end{eg}

特定の位数の正規部分群が無いことを示すには, 共役類の直和を考えることが一つの方法としてある. 

\begin{egprob} 
$D_6$は位数4の正規部分群を持たないことを示せ. 
\end{egprob}
\begin{point}
正規部分群は共役類の直和なので, 共役類をすべて調べてしまって, その直和が位数4の群を作り得ないことを示せばよい. 共役元の計算は, $\tau\sigma\tau = \sigma^{-1}$などの関係式を使って効率的に計算する. $D_6$は回転と鏡映を表す元で生成できることにも注意. 
\end{point}
\begin{egans}
$D_6$には回転(位数6の元)$\sigma$と, ある軸による鏡映$\tau$(位数2)が存在する. $D_6$の任意の元$g$が$g= \sigma^k \tau^i$ ($k=0,\dots, 5, i=0,1$)の形に一意に書ける. $\tau \sigma \tau^{-1} = \sigma^{-1}$や, それと同値な$\sigma\tau = \tau\sigma^{-1}$に注意すると
\[\tau g\tau^{-1} = \tau\sigma^k \tau^{i-1} = \tau\sigma^k\tau^{-1}\tau^i = \sigma^{-k}\tau^{i}\]
\[\sigma g\sigma^{-1} = \sigma^{k+1}\tau^{i}\sigma^{-1} = \bcas  
\sigma^{k} \quad (i=0) \\
\sigma^{k+2}\tau\quad (i=1)
\ecas\]
が得られる. ここから, $g=\sigma^k\tau^i$の共役元をすべて計算することができる. 
\[\sigma^m g\sigma^{-m} = \bcas 
\sigma^k \quad (i=0) \\
\sigma^{k+2m} \quad (i=1)
\ecas \]
\[\sigma^m \tau g \tau^{-1} \sigma^{-m} = 
\bcas 
\sigma^{-k} \quad (i=0)\\
\sigma^m \sigma^{-k}\tau \sigma^{-m} = \sigma^{-k+2m}\tau\quad (i=1)
\ecas
\]
となる. 上を見れば, 共役類を解析することができ, 以下が$D_6$の共役類である. 
\ben[label=$\bullet$] 
\item $R_0=\set{e}$ ($e$は単位元)
\item $R_1=\set{\sigma, \sigma^5}$ 
\item $R_2=\set{\sigma^2, \sigma^{4}}$
\item $R_3=\set{\sigma^3}$ 
\item $S_0=\set{\sigma\tau,\sigma^3\tau, \sigma^5\tau}$ 
\item $S_1=\set{\tau,\sigma^2\tau, \sigma^4\tau}$ 
\een 
正規部分群は, その定義から, 共役類の直和にならねばならない. 上の共役類の合併を考えて位数4の群が作れるか, というパズルを考えればよい. そしてそれは不可能であることがわかる. 仮に位数4の正規部分群$N$が作れたとしよう. まず$N$は単位元$e$を含む. $N$は$R_i,S_i$の合併でかつ位数4である. 位数4であることから考えられる組み合わせは, $N=\set{e}\sqcup S_0, \set{e}\sqcup S_1, \set{e}\sqcup R_3\sqcup R_1, \set{e}\sqcup R_3\sqcup R_2$などが考えられよう. しかし, これらはいずれも群をなさないので適さない. (たとえば最初の候補では, $\sigma^2\tau \cdot \tau = \sigma^2$が$\set{e}\sqcup S_0$に属していないのでおかしい. 最後の候補では, $\sigma^2\cdot \sigma^3 = \sigma^5$が合併に属さない.) ゆえに, これら共役類の直和によって位数4の群を作ることは不可能であり, 位数4の正規部分群は存在しない. 

\end{egans}



\begin{egprob} (H28東工大, 午後1) 
有限群$G$の位数を$n$とし, $G$の共役類全体を$C(1),\dots, C(r)$とする. 
\ben 
\item $|C(i)| = n(i)$とすると$n(i)$は$n$の約数であることを示せ. 
\item $l(i) = n/n(i)$とおくとき$\sum_{i=1}^{r} 1/l(i) = 1$を示せ. 
\item $r=1,2,3$となる$G$の位数$n$をすべて求めよ. 
\een 
\end{egprob}
\begin{egans}
(1): $x\in C(i)$とする. $C(i)$には$G$が共役で作用する. $x$の固定部分群$G_x$は$C_G(x)$である. よって, $|C_G(x)||C(i)| = |G|$だから$n(i) = [G:C_G(x)]$. \\
(2): 類等式から明らか. \\
(3): (2)から考えると$n=1,2,3,4,6$に限られ, $n=4$はアーベル群で共役類が4つになるので不適. ほかは全部よい($n=6$のときは$S_3$である). 
\end{egans}





\subsection{Abel群の構造定理}
\begin{theo} (有限Abel群の構造定理)
\end{theo}
\begin{theo} (有限生成Abel群の構造定理)

\end{theo}


\begin{eg}
$p$を素数とする。位数$p^2$の元はAbel群であった。よって, 位数$p^2$の群は$\mb{Z}/p^2\mb{Z}$か$\mb{Z}/p\mb{Z}\times \mb{Z}/p\mb{Z}$に同型である。
\end{eg}

\begin{egprob} (H15数学I)\\
自然な準同型
\[\phi:GL_{2}(\mathbb{Z}/4\mathbb{Z}) \to GL_{2}(\mathbb{Z}/2\mathbb{Z})\]
の核は$(\mathbb{Z}/2\mathbb{Z})^4$と群として同型であることを示せ。
\end{egprob}
\begin{egans}
$a,b,c,d\in \mb{Z}$として, これを$\mb{Z}/n\mb{Z}$の元と区別せずに用いる。
$A=\bep a & b \\ c & d \enp \in \ker{\phi}$とする。このとき, 
\[a,d\equiv 1\pmod{2},\quad  c,b\equiv 0\pmod{2}\]
逆にこのような$a,b,c,d$が与えられると$\det{A} = ad-bc$は奇数で代表されるので, $\mb{Z}/4\mb{Z}$の単元である。よって$A\in GL_2(\mb{Z}/4\mb{Z})$である\footnote{可換環$R$で, $X\in M_n(R)$に対し$X\in GL_n(R)\iff \det{X} \in R^{\times}$である。余因子行列は代数的に作れるので線形代数とほとんど同じである。}。よって$\ker{\phi}$の位数が$16$であることが分かる。\par 
あとは次を示せばよい。
\ben 
\item $\ker{\phi}$のすべての元の位数が$2$以下であること。
\item $\ker{\phi}$がAbel群であること。
\item $(\mb{Z}/2\mb{Z})^4$以外の可能性がないこと。
\een
これらは簡単な整数問題になるので, そこは読者にお任せする。
\end{egans}
\begin{prac}
上の3つを示し, 証明を完成させよ。
\end{prac}





\newpage 
\subsection{Sylowの定理}
かなり強力な群論の定理のひとつである。
\tbox{ 
\begin{theo} (\tf{Sylowの定理}) \label{theo:sylow} 
$G$を有限群, $p$を素数, $|G| = p^{a} m$ ($a>0$, $m$は$p$で割れない自然数)とする。このとき, 次の(1)-(4)が成り立つ。
\ben
\item $G$の部分群$H$で$|H| = p^{a}$となるものが存在する(このような部分群をSylow$p$部分群という)。
\item Sylow$p$部分群$H$をひとつ固定する。$K\subset G$が部分群で$|K|$が$p$冪なら, $g\in G$があり$K\subset gHg^{-1}$となる。特に, $K$を含む$G$のSylow$p$部分群が存在する。
\item \tf{$G$のすべてのSylow$p$部分群は共役である.}
\item Sylow-$p$部分群の数$n_p$は
\[\bolm{n_p = |G|/|N_{G}(H)| \equiv 1 \pmod{p}}\]
という条件を満たす.
\een
\end{theo}
}{title=Sylowの定理}

\begin{cor} (\tf{Cauchyの定理})\\
$G$を有限群, $p$を$G$の位数を割り切る素数とするとき, 位数$p$の元$g\in G$が少なくとも一つ存在する. 
\end{cor}
\prf Sylow-$p$部分群$P$を取り, $a\in P$は$G$の単位元と異なるとする. このとき, $a$の位数は$p^{n}$の形で書けるので, $a^{p^{n-1}}$は位数$p$の元である. \qed


\begin{egprob}
$D_6$には位数4の正規部分群は存在しないことを示せ. 
\end{egprob}
\begin{egans}
$D_6$のSylow-3部分群は正規である. [証明]もし正規でないなら$4$個あるはずである. このとき位数3の元が8個見つかるが, $D_6$には位数2の元が$6$個はあるので, おかしい. よってSylow-3部分群$P_3$が一つあり, 正規.\par 次にSylow-2部分群は正規ではない. 実際, もし正規だとしてそれを$P_2$とすると, $D_6 = P_3P_2$でかつ$P_2\cap P_3 = \set{e}$かつ$P_2,P_3$が正規なので, $f:P_2\times P_3 \to D_6$; $f(x,y) = xy$が群同型となるので,$D_6\cong P_2\times P_3$になるが, 右辺はAbel群なので矛盾. ゆえに位数4の部分群は正規ではありえない. \qed 
\end{egans}



\begin{eg} (\tf{位数$pq$の群; $p<q, p\not| q-1$})

\end{eg}

\begin{egprob} (都立大\footnote{とはいえ, これは昔の問題である.  最近の院試でこの手の出題は見ないのであんまり恐れなくてもいいかもしれない. })
位数$p^2q$の群の型を決定せよ. ($p<q$は素数, $p\not| q-1$)
\end{egprob}

\begin{egans}
条件より$n_p = 1$. よって$S_p\subset G$は正規. $n_q$の可能性は$p<q$にも注意すると$n_q \in \set{1,p^2}$. $n_q = 1$ならば$S_q$も正規で$G\cong S_p\times S_q$で, $C_{p}^2 \times C_q$もしくは$C_{p^2}\times C_q$に同型である. 以下$n_q = p^2$で考える. $q|p^2-1 = (p-1)(p+1)$で$p-1<q$だから$q|p+1$であり, $p<q\leq p+1$なので$q=p+1$が従う. $p,q$は素数だから$p=2,q=3$となるが, これは$p\not| q-1$に反する. 以上ですべて決定できた. \qed 
\end{egans}

\begin{eg}
位数21の群で非可換なものがただひとつ存在することが知られている. それは次のアファイン変換群の形で実現できる. 
\[G=\set{\phi_{\alpha,\beta}\in \mr{Aut}(\mb{F}_{7})\mid  \phi_{\alpha,\beta}(x) = \alpha x + \beta,\quad \alpha,\beta\in \mb{F}_{7},\alpha^3=1}\]
なお, このアファイン変換群の出どころはこの次の節の\tf{半直積}と呼ばれるものから構成できる. 興味がある人は見てみよう. 
\end{eg}

ガロア理論と絡めた問題を考えよう. 
\begin{egprob}  
$\mb{Q}(\sqrt[11]{2}, \zeta_{11})/\mb{Q}$が110次Galois拡大であることを示せ. また, 中間体$M$で$\mb{Q}$上55次, 22次, 10次のものはいくつあるかそれぞれ求めよ. 
\end{egprob}
\begin{egans}
Galois拡大であることは省略($11$が素数であることは使う). $110 = 2\times 5\times 11$である. 55次のものは位数2だからSylow-2部分群に対応する. $G$をガロア群とする. \ao{Sylow部分群はどの二つも共役なので, 55次中間体$M$を一つ見つけると$\sigma(M)$ですべて得られる(ここがポイント).} $M=\mb{Q}(\sqrt[11]{2},\zeta_{11} + \zeta_{11}^{-1})$であり, $\sigma(M)$はある整数$j$によって$\mb{Q}(\sqrt[11]{2}\zeta_{11}^{j}, \zeta_{11}+\zeta_{11}^{-1})$の形になることが分かる. これらは異なる$j\in \cyc{11}$で異なる体になることが示せるので, 11個である. 22次の場合, $M=\mb{Q}(\sqrt[11]{2}, \zeta_{11} + \zeta_{11}^{4} + \zeta_{11}^{5} + \zeta_{11}^{9} + \zeta_{11}^{3})$は一つの例であり, 11個あることが示せる. 10次の場合, 実は$\mb{Q}(\zeta_{11})$しかない. 
\end{egans}
\begin{prac}
上の解答の細部を埋めよ.  
\end{prac}










\newpage 
\subsection{($\spadesuit$)半直積} 
ガロア理論の院試で一回だけ半直積を見たことがあり「ふざけんな」となったのでまとめた. 半直積とは直積群の一般化である. 集合としては直積に等しいが, その群構造が少し変わっている. 二つの群から新たな群を作り出す方法であり, 有限群の分類にもしばしば現れる. 
\tbox{ 
\begin{defn}
群$N,H$と群準同型$\phi:H\to \mr{Aut}(N)$を与える. $\phi(h)$のことを$\phi_{h}$で表す. このとき, $\phi$に関する$N$と$H$の\tf{(外部)半直積}$N\rtimes_{\phi} H$とは, 次のようにして定まる群である. 
\ben[label=$\bullet$] 
\item 集合としては$N\rtimes_{\phi}H = N\times H$. 
\item \bolm{(n_1,h_1)(n_2,h_2) = (n_1\phi_{h_1}(n_2), h_1h_2)}
\een 
\end{defn}
}{title=半直積}

\begin{prac}
$N\rtimes_{\phi}H$が群をなすことを確認せよ. 単位元は$(1_N,1_H)$で, $(n,h)$の逆元が$(\phi_{h^{-1}}(n^{-1}), h^{-1})$であることも確認せよ. 
\end{prac}

\begin{prac}
この演算はいかにも非可換っぽく見えるから, うまくやると二つの有限群から非アーベル群が作れるということが予想される. \ao{素数$p,q$で$q-1$が$p$を割り切るとき, 位数$pq$の非アーベル群が存在するが, それは実は半直積で構成されるのである.} だが$p,q$に関する条件が満たされないときはアーベル群しか存在し得なかった. どうしてこのような場合に位数$pq$のアーベル群が半直積から構成できないのだろうか. 
\end{prac}
%Answer: Hは位数p, Nは位数qとする. \phi: H\to Aut(N)\cong (\cyc{q})^{\times} という準同型は位数pから位数q-1の射になる. pは素数なのでこの射は単射か自明な準同型だが, 単射ならばq-1がpの約数にならねばならないはず. よって\phiが自明な準同型しかないので, \phi_{h_1}は常に恒等射N \to Nになる. すると, 半直積の演算は直積群の演算になるので, (HとNは素数位数でアーベル群だから)非アーベルにならない. 


\begin{eg} 個人的には次の例が最も良い半直積の例である. $\mb{R}^2$のアファイン変換$f_{a,b}(x) = ax+b$ ($a\in \mb{R}^{\times}, b\in \mb{R})$の集合$G$は合成に関して\tf{非可換な}群をなす. 実際, 実数$a_i\neq 0,b_i$ ($i=1,2$)に対して
\bealn{ 
f_{a_1,b_1}\circ f_{a_2,b_2}(x) &= a_1(a_2x+b_2) + b_1 \\
&= a_1a_2x + (a_1b_2 + b_1) = f_{a_1a_2, a_1b_2 + b_1}(x)
}
であるから. そして, この結果を見るとまさしく半直積が現れている. 実際, この例では$G\cong \mb{R}\rtimes_{\phi} \mb{R}^{\times}$である. ただし, $\phi:\mb{R}^{\times}\to \mr{Aut}(\mb{R})$は$t\in \mb{R}^{\times}$に対し$\phi_{t}(r) = tr$なる$\phi_t$を与える写像である. $N=\mb{R}$のほうは加法群であることに注意. また, この$\mb{R}$は一般の体に置き換えてもよい(環まで許すと$t$倍が自己同型とは限らないので作れなくなる). 
\end{eg}

\begin{eg}
Sylowの定理の節で考察した群は$\mb{F}_{7}\rtimes_{\phi} \mb{F}_{7}^{\times}$である. ただし$\phi: \mb{F}_{7}^{\times} \to \mr{Aut}(\mb{F}_{7})$は上の例と同様の, $\mb{F}_{7}$の元の掛け算による自己同型である. 
\end{eg}

外部半直積は独立した群から一つの群を得る操作だが, 一つの群に内包された状況における\tf{内部半直積}というものもある. 
\begin{defn}
ふたつの群$N,H$に対して, $N$の$H$による\tf{内部半直積}とは, 次の性質を満たす群$G$のことで, $G=N\rtimes H$と表す. 
\enu{ 
\item $N$は$G$の正規部分群で$G=NH$ 
\item $N\cap H = \set{1}$
}{}
\end{defn}




\tbox{ 
\begin{prop} (参考:Wikipedia「半直積」)
\ben 
\item 標準的な単射$N\hookrightarrow N\rtimes_{\phi} H$, $H\hookrightarrow N\rtimes_{\phi}H$が存在する. 
\een 
\end{prop}
}{title=半直積のさまざまな性質}

$G=NH, N\cap H = \set{1}$という状況は有限群論でよく現れるが, この条件は次のように特徴づけられる. ようは内部半直積とは, $N$と$H$の元による積で分解する操作が一意であるという状況に現れるのである. 
\tbox{ 
\begin{prop} (参考: Wikipedia「半直積」)
$G$の正規部分群$N$と部分群$H$に対し, 次は同値. \ben 
\item $G=NH$かつ$N\cap H = \set{1}$
\item $G$のすべての元は積$nh$ ($n\in N, h/in H$)として一意的に書ける. 
\item $G$のすべての元は積$hn$ ($n\in N, h/in H$)として一意的に書ける.
\een 
\end{prop}
}{}



次のような状況がしばしばあり, 有用である. 
\tbox{ 
\begin{theo} (参考:Wikipedia「半直積」)
群$G$と正規部分群$N$と部分群$H$を与える. 
$G$のすべての元が$g=nh$ ($n\in N, h\in H$)の形に書けるとする. $\phi:H\to \mr{Aut}(N)$を$\phi_h(n) = hnh^{-1}$で与える(これは群準同型である). このとき\bolm{G\cong N\rtimes_{\phi} H}である. ただし, 同型写像は$nh$を$(n,h)$に送る. 
\end{theo}
}{title=半直積に同型であるための十分条件}
 
 \begin{egprob} $p,q$は素数, $p|q-1$とする. 位数$pq$の非アーベル群$G$は同型除いてただひとつであることを証明せよ. 
 \end{egprob}
\begin{egans}
$G$から位数$p,q$の元が取れる. それで生成される巡回部分群を$H,N$とする. $p<q$なので, Sylow-$q$部分群である$N$は正規である. $\phi:H\to \mr{Aut}(N)$を$\phi_h(n) = hnh^{-1}$で定める. $\phi$は単射か自明な準同型であるが, 自明な準同型だと非アーベルに反するので単射.  とくに$H$の生成元が$\sigma$なら$\phi_\sigma$の位数が$p$であるから, $\phi_{\sigma}^{p} = \mr{id}_N$により$\sigma^{p} n \sigma^{-p} = n$を得る. これが任意の$n\in N$で成り立つことは, $N$の生成元$\tau$に対して$\sigma^{p} \tau \sigma^{-p} = \tau$であるということに同値である. $G$の元は$nh$ ($n\in N, h\in H$)の形で書ける($G=NH, N\cap H = \set{e}$より). よって
\[G = \lan \sigma,\tau\mid \sigma^p = \tau^{q} = 1, \sigma^p\tau\sigma^{-p} = \tau \ran\]
が得られる. 
\end{egans} 
 
 
 
 半直積でお手軽に非アーベル群が作れるものだから, 院試の題材にはもってこいだろう. 
 
 
次はパズルチックだが良問である. 位数21非アーベル群の例も習得できる. 
\begin{egprob} (R2 東北大 専門 1)
集合$\mb{Z}/7\mb{Z}\times \mb{Z}/3\mb{Z}$に次の積を定めた群を$G$とする. 
\[(a,b)(x,y) = (a+2^x b, x+y)\]
\ben 
\item $H = \set{(a,0) \mid a\in \mb{Z}/7\mb{Z}}$, $K = \set{(0,x)\mid x\in \mb{Z}/3\mb{Z}}$はそれぞれ正規部分群か?
\item $G$の中心を求めよ. 
\item $G$の共役類の個数を求めよ. 
\item 位数が3である$G$の元の個数を求めよ. 
\een
\end{egprob}

この群は$\cyc{7}\rtimes_{\phi}\cyc{3}$で, $\phi:\cyc{3}\to \mr{Aut}((\cyc{7})^{\times})$は$\phi_{x}(b) = 2^a b$に取っているということである. 
\begin{egans}
(1): Sylowの定理より$n_7 = 1$なので$H$は正規. もし正規なら$n_3 = 1$だが, このとき$HK$は部分群で$G = H\times K$となるので非アーベルに反する. よって正規ではない. \\
(2): 計算するのみ. $Z(G) = \set{1}$. \\
(3): 類等式. $C$は2個以上の元を持つ共役類を走るとき
\[21 = |Z(G)| + \sum_{|C| \geq 2} |C|\]
である. $|C|$は21の真の約数なので3か7. $21-1 = 20$を3と7の足し算で分解する方法は$3+3+7+7$なので, 共役類は($\set{1}$と合わせると)5個. \\
(4): 共役類を$\set{1}, C_{3,1}, C_{3,2}, C_{7,1},C_{7,2}$と名付ける, $|C_{3,i}| = 3$, $|C_{7,i}| = 7$である. \\
$H$が正規部分群なので, 共役類のdisjoint unionになることから組み合わせを考えると$H=\set{1}\sqcup C_{3,1}\sqcup C_{3,2}$となる. $(a,x)^{-1} = (b,-x)$と書けるが\aka{(\tf{注意}: $b=-a$ではない！)}, 
$(a,x)(0,1)(b,-x) = (\ast, 1)$より分かるように$(0,1), (0,2)$は共役ではないので同じ共役類に属さず, 位数3の元なので$(0,1), (0,2)\notin H$より, $(0,1)\in C_{7,1}, (0,2)\in C_{7,2}$としてよい. ひとつの共役類上で位数が不変量であるから$C_{3,i}$のすべての元は位数7, $C_{7,i}$のすべての元は位数3である. よって, 14個である. 

\end{egans}
\begin{prac}
上の問題について, (1)で$K$が正規でないことを計算で確かめよ. (2)を確かめよ. 正規部分群が共役類の合併であることを一般に証明せよ. 
\end{prac}
 
 
 
 
 
 
 
\subsection{その他}
問題で使った知識を置いておく. 
\begin{prop}
$\mb{Z}/n\mb{Z}$に対し, $\mr{Aut}(\mb{Z}/n\mb{Z}) \cong (\mb{Z}/n\mb{Z})^{\times}$.
\end{prop}
\prf $\phi:\mb{Z}/n\mb{Z} \to \mb{Z}/n\mb{Z}$は$\phi(1)$の値で決まる. $\phi(1)$が$n$と互いに素な類にならないと全射にならない. 逆に$n$と互いに素の類になれば$\phi(1)$が$\mb{Z}/n\mb{Z}$で単元だから$\phi(1)^{-1}$倍の自己同形も考えることができて, それによって全単射になる. ,
\qed 



\subsection{例題集}
\begin{egprob} (2020前期代数演義)
位数595の群は位数17の正規部分群を持つことを示せ。
\end{egprob}
\begin{egans}
$595=5\times 7\times 17$より$n_{17}$としてありえるのは$1$か$35$である。$n_{17} \neq 35$を示せばよい。$n_{17} = 35$だとする。35個ある位数17の部分群は 単位元でのみ交わるので, $G$に位数17の元が$n_{17}\times (17-1) = 560$個以上あることが保証される。
いま, $n_{17} = 35$のもとでは$n_{7} = 1$となることを示す。
実際, $n_{7}$としてありうるのは1か85であり, もし85であれば位数$7$の元が$85\times (7-1)$個以上あることになり, 位数7,17の元だけで$595$個以上の元が保証されてしまうので矛盾である。
よって$n_{7} = 1$となる。$S_{p}$をsylow-$p$-部分群とすると, $S_7 \cdot S_{17} < G$で, $S_7\cap S_{17}= \{ 1 \}$だから$|S_{7}\cdot S_{17}| = 119$となる。
$S_{7}\cdot S_{17}$の中でSylowの定理を用いて, $17\not\equiv 1 \pmod{7}$, $7\not\equiv 1\pmod{17}$から
$S_{7}, S_{17} \rotatebox{90}{$\triangle$} S_{7}\cdot S_{17}$である。
$S_{7}\cap S_{17} = \{ 1 \}$だから$S_{7}\cdot S_{17} \cong S_{7}\times S_{17}$である。よって, $G$ には位数$7\times 17$の元が
\[7\times 17 - 7 - 17 + 1 = 96\]
個あることが保証されるが, すでに位数17の元が$560$個保証されているので矛盾である。よって$n_{17} = 1$であり, 定理\ref{theo:sylow}(3)により正規部分群である。\qed
\end{egans}

\begin{egprob} 
$n,a$を自然数とする。$\phi(a^{n} - 1)$は$n$で割り切れることを示せ。
\end{egprob}
\begin{egans}
$G = (\mb{Z}/(a^{n} - 1)\mb{Z})^{\ast}$ とする。$a$と$a^{n} - 1$は互いに素なので$a\bmod{a^{n} - 1}$は$G$に属す。$a \bmod{a^{n}-1}$の位数が$n$であることは易しい。よって$n$は$|G|$を割り切る。$|G| = \phi(a^{n} - 1)$だから主張が示された。
\end{egans}


\clearpage
\subsection{演習問題Part1}
\subsubsection{Level A}
\begin{prob} (9/26, Quiz)
$G,H$は有限群であり,$|G|=n,|H|=m$としたとき$n$と$m$は互いに素とする.このとき,準同型写像$f:G\to H$を全て求めよ.
\end{prob}

\begin{prob} (9/20, Quiz)
$G$を位数$n$の有限群とするとき, $\text{Hom}(G, \mathbb{C}^{\times})=\text{Hom}(G, \langle \zeta_n\rangle )$を示せ. ただし$\zeta_n$は$1$の原始$n$乗根で, $\text{Hom}(G, H)$は群$G$から群$H$への準同型全体. (さらに$G$が巡回群なら$|\text{Hom}(G, \mathbb{C}^{\times})|=n$となる.)
\end{prob}

\begin{prob} (2020前期代数演義)
$X=\bep 0 & 1 \\ 1 & 0 \enp \in GL(2,\mb{R})$とする。
\[i_X(A) = XAX\quad (A\in SL(2,\mb{R}))\]
と定義するとき, $i_{X}$は外部自己同型であることを示せ。つまり, 内部自己同型ではないことを示せ。
\end{prob}


\begin{prob} (H24数学I-2)
$p\geq 3$を奇素数, $n$を自然数とする. 
\[G = \set{\bmat{a&b\\ 0&d} \mid a,b,d\in \cyc{p^n}, ad=1}\]
には位数$p^{2n-1}$の部分群がただひとつ存在することを示せ. 
\end{prob}


\subsubsection{Level B}


\begin{prob} (H28東北大選択1)
正の整数$m,n$に対して$m$次対称群$S_m$の位数$n$の部分群の個数を$N(m,n)$とおく。
\ben
\item $N(5,3)$を求めよ。
\item $N(5,4)$を求めよ。
\item 整数$m\geq 2$に対し, 次の式を示せ。
\[N(m,2) = \disp\sum_{k=1}^{[\frac{m}{2}]} \dfrac{m!}{2^{k}k!(m-2k)!}\]
\een
\end{prob}

\subsubsection{Level C}
\begin{prob} (2020前期代数演義, 星問題)
$SL_{2}(\mb{F}_{3})$のSylow-2-部分群は正規でないことを証明せよ。
\end{prob}

\begin{prob} (H30専門科目2) \label{prob:H30senmon-2}
$p$は素数, $k,m$を正の整数で, $k$と$p^2-p$は互いに素であるとする. 位数$kp^m$の有限群$G$が次の性質を満たす部分群$N,H$を持つとする. 
\ben 
\item $N$は位数$p^m$の巡回群で$G$の正規部分群である. 
\item $H$は位数$k$の群である. 
\een 
このとき$G$は$N$と$H$の直積であることを示せ.
\end{prob}

\begin{prob} (H27専門科目)\label{prob:H27senmon-1}
$G$は非可換群で次の条件$(\ast)$を満たすとする. 
\begin{center}
    $N_1,N_2\subset G$が相異なる自明でない($\set{1},G$と異なる)正規部分群なら$N_1\not\subset N_2$である. 
\end{center}
\ben 
\item $N_1,N_2$が相異なる$G$の自明でない正規部分群なら$G=N_1\times N_2$であることを証明せよ. 
\item $G$の自明でない正規部分群の可zは高々2個であることを示せ. 

\een 
\end{prob}






\subsubsection{Hint・Answer}
{\small 

\begin{probans}

\end{probans}

\begin{probans}

\end{probans}

\begin{probans}

\end{probans}

\begin{probans}
位数$p^{2n-1}(p-1)$で, Sylow$p$部分群の個数$n$は$n-1\equiv 0\pmod{p}$かつ$p-1$の約数. よって$n=1$なので示せた. 

\fbox{別解}$(\cyc{p^{n-1}})^{\times}$が巡回群であることを用いよ($p=2$を除外するのはこのため). $a,d$はこの乗法群に属す. この乗法群の指数$p-1$の部分群$A$を取り, $G$の元で対角成分を$A$に制限した元全体のなす部分群を$H$とすると$H$は位数$p^{2n-1}$の部分群である. これはSylow-$p$部分群であるが, Sylow-$p$部分群はすべて共役のはずである. しかし$H$が正規部分群であることは容易に確かめられるので$H$の一つに限られる. 

\end{probans}
%=============================Level B



\begin{probans}

\end{probans}

\begin{probans}

\end{probans}

\begin{probans}
$\phi:H\to \mr{Aut}(N)$について調べよ. $h\in H$, $n\in N$について $hnh^{-1} = n$を示せ. 
\end{probans}

\begin{probans}
(i): $N_1N_2$と$N_1\cap N_2$が正規. よって
$N_1\cap N_2 = \set{e}$が分かる. ゆえに$N_1\times N_2 \to G$; $(n_1,n_2)\mapsto n_1n_2$が群同型($N_1,N_2$の元が互いに可換であることを用いると準同型であることが言える). \\
(ii): $N_1,N_2,N_3$が相異なる非自明な正規部分群とする. このように3つ取れると, $n_2,m_2\in N_2$に対し, $n_2m_2 = m_2n_2$が成り立つ($n_2 = a_1a_3\in N_1N_3$と書けばわかる). $N_1,N_2$の元は可換で, $G=N_1N_2$だから,  $G$はアーベル群になり, 矛盾. なお, $G$がアーベル群の場合に条件を満たすが正規部分群がいくつもある例として, $(\mb{Z}/p\mb{Z})^2$がある. 
\end{probans}
}






\newpage 
\section{群準同型}
\begin{prop}
$\phi:G\to H$を群準同型とする。
\ben 
\item \tf{$\ker{\phi}$は$G$の正規部分群である。}
\item $G/\ker{\phi} \cong \im{\phi}$
\een
\end{prop}


\begin{egprob} (2013基礎数学試験No.2 一般化)
$S_n$からアーベル群$G$の全射群準同型$\phi$が存在するならば, $G=\{ 1 \}, \mb{Z}/2\mb{Z}$であることを示せ。
\end{egprob}
\begin{egans}

$|G| \geq 3$を仮定する。 $G$がアーベル群であるという仮定によって$\phi(hgh^{-1}) = \phi(g)$が成り立つことに注意する。よって, $\phi$は共役類上では一定の値を取る。すべての互換は共役であり, ある互換が$\phi$で単位元に移ると, $\phi$は自明な準同型になるので不適。よって$\phi$はすべての互換を, $G$のある位数2の元$g_2$に移す。$S_n$は互換で生成されるので, $\phi$の像は$\{ 1 , g_2\}$に含まれる。これは2元集合だから, 全射になり得ない。また, $|G| = 2$ならば符号を取る準同型$\mr{sgn}$は$\mb{Z}/2\mb{Z}$への全射準同型である。\qed
\end{egans}

\subsection{準同型定理}

\begin{egprob} (H26-基礎科目II-4)
$G = \mb{Z}/4\mb{Z}\times \mb{Z}/6\mb{Z}\times \mb{Z}/9\mb{Z}$の指数3の部分群は何個あるか. 
\end{egprob}
\begin{egans}
\tf{ $\phi:G\to \mb{Z}/3\mb{Z}$という全射群準同型をすべて求めて, $\ker{\phi}$として挙がるものをすべて求める. } なぜこれでよいかというと, 任意の指数$3$の部分群$H$に対し$\psi :G/H \to \mb{Z}/3\mb{Z}$という同型が一つあり, これによって$G\to G/H \to \mb{Z}/3\mb{Z}$は$H$を核に持つ全射準同型になるからである. また, 中国剰余定理より
\[G\cong \mb{Z}/4\mb{Z}\times \mb{Z}/2\mb{Z} \times \mb{Z}/3\mb{Z} \times \mb{Z}/9\mb{Z}\]
である. $G$の元を$(x,y,z,w)$で表すとき, $\phi(1,0,0,0) = \phi(0,1,0,0) = 0$が分かる(Why?). $\phi(0,0,1,0), \phi(0,0,0,1)\in \mb{Z}/3\mb{Z}$で$\phi$は一意に決まり, 全射性よりどちらかは0ではない. このことから8個の$\phi$があると分かり, それらから$\ker{\phi}$をすべて計算すると4個異なる部分群が出てくる. よって4個である. 
\end{egans}
\begin{egans} (T氏のアイデアをもとにしたより良い解答)\\
$H$を指数$3$の部分群とする. $G/H \cong \mb{Z}/3\mb{Z}$なので\ao{$3G  \subset H$である.(ここがポイント)}  よって$H$は$G/3G \cong \mb{Z}/3\mb{Z} \times \mb{Z}/3\mb{Z}$の中の指数$3$の部分群$H/3G$に対応する. それは4個ある. 実際, 指数3ということは位数3なので, 位数3の巡回元を探せばよく $(1,0)$, $(0,1)$, $(1,1)$, $(1,-1)$で生成されるものがそれらである. 
\end{egans}


\begin{prob} (Liの自作問題, 難)
$G = \prod_{n=1}^{2021} \mb{Z}/n!\mb{Z}$とする. $G$の指数$2017$の部分群はいくつあるか. (2017は素数である)
\end{prob}

\begin{egprob}
群$G$が次で与えられるとき, 群準同型$\phi:G\to \mathbb{C}_{1}$をすべて求めよ\footnote{($\spadesuit$)表現論の話。$n$次元複素ベクトル空間$V$と有限群$G$に対し, 「群$G$の$n$次元(既約)表現」というのは群準同型$G\to \to GL(V)$のことであり, とくに$n=1$のときは$GL(V)\cong \mb{C}_{1}$だから, この問題は「1次元既約表現の数を決定せよ」という問題に他ならない。}。ただし, $\mb{C}_{1}$は絶対値1の複素数の集合に, 通常の積による演算を入れた群である。
\ben 
\item $\mb{Z}/n\mb{Z}$
\item Kleinの四元群$K = \{ 1,(12)(34),(13)(24),(14)(23) \} \subset \mathfrak{S}_{4}$
\item $\mb{Z}/2\mb{Z}\times \mb{Z}/2\mb{Z}$
\item $\mb{Z}/2\mb{Z}\times \mb{Z}/4\mb{Z}$
\item $\maf{S}_n$ (\tf{Hint:}互換で生成できることを用いよ。)
\item (なんでも考えてみてください！)
\een
\end{egprob}
\begin{egans}
(1): $\phi(1)$の行先は1の$n$乗根であり, その行先で準同型は決定されるので$n$個。\\
(2) 
\end{egans}



\subsection{群の同型と非同型}
簡単な群の同型の例を紹介するとともに, 群論的な議論を総合的に復習する。

\begin{eg} ($PSL_{2}(\mb{F}_{2}) \cong \maf{S}_{3}$, 青雪江Ex2.6.1)\\
$G=PSL_{2}(\mb{F}_{2})$とする。$G$は$SL_{2}(\mb{F}_2)$の元で定数倍で一致するものをすべて同一視することで得られるが, 各同値類は1つの元しか持たないので$SL_{2}(\mb{F}_{2}) \cong G$である。$SL_{2}(\mb{F}_{2})$の単位元ではない元は,
\beqn
\bep
1 & 1 \\
0 & 1
\enp
,
\bep
1 & 0\\
1 & 1
\enp
,
\bep
0 & 1\\
1 & 0
\enp 
,
\bep 
0 & 1\\
1 & 1 
\enp
,
\bep 
1 & 0\\
1 & 1
\enp 
\eeqn
である。このうち最初の3つだけが位数が$2$である。$X$をこの3つの元の集合とする。共役作用によって$G$は$X$に作用し, したがって準同型$\phi:G\to \maf{S}_{3}$を得る。すなわち, 
\[\phi(A) = (X\ni M\mapsto A^{-1}MA\in X)\]
とおくのである。$\phi$は同型$G\cong \maf{S}_{3}$を定める。これを見るには$\phi$が単射であることを見ればよいので, すべての$M\in X$で$MA = AM$であるならば$A = 1_{G}$であることを示せばよい。これは簡単に示せる。
\end{eg}

位数が同じだが同型でないことが分かる例を挙げる. しばしば, 同型でないことを論じるには\tf{元の位数}が活躍する. 

\begin{egprob}
次の群は互いに同型ではないことを示せ。ただし, $\mb{Z}_{d} = \mb{Z}/d\mb{Z}$
\ben 
\item $\mb{Z}_2\times \mb{Z}_2$と$\mb{Z}_{4}$
\item $S_4$と$SL_2(\mb{F}_{3})$
\een
\end{egprob}














\section{群の作用}


\subsection{軌道}
群は作用を考えてからが本番。
\begin{prop}
有限群$G$が集合$X$に作用しているとする. $x\in X$の固定部分群を$G_x$, 軌道を$O_x$で表すとき, 
\ben 
\item \bolm{|G_x||O_x| = |G|}        
\item 
\een
\end{prop}
\prf (1): $a\in O_x$とする. $O_x = \set{x,g_2x, \dots, g_n x}$とし, これらに重複がないとする($|O_x| = n$). 写像$\phi:G\to O_x$を$\phi(h) = hx$により定めると, $hx = g_i x\iff g_i^{-1}h \in  G_x$である. よって, 任意の$i$に対して $\phi^{-1}(g_ix) = g_i G_x$だから
\[G = \bigcup_{i=1}^{n} \phi^{-1}(g_i x) = \bigcup_{i=1}^{n}g_iG_x \]
から集合の元の個数を計算すれば得られる. \qed 

\begin{eg}
位数$2m$ ($m$は奇数) の群$G$において, 位数2の元はSylow-2部分群を作る. Sylowの定理から位数2の元はすべて共役であることが分かる. よって, 位数2の元全体は一つの共役類$X$をなす(共役類上で位数は一定であることに注意). $X$には$G$が共役で作用する.  もしある位数2の元$x$が$\cyc{2k}$ に同型な部分群に含まれているとする. このとき, $x$の固定部分群$G_x$, すなわち, すなわち中心化部分群$C_{G}(x)$はその部分群を含んでいるので, $|X|$は$2k$の倍数であることが分かる. 
\end{eg}







\subsection{cosetへの作用}

次の問題は覚えておくとよい。
\tbox{ 
\begin{egprob} 
$G$を有限群, $p$を$G$の位数を割り切る最小の素数とする。任意の指数$p$の$G$の部分群は正規部分群であることを証明せよ。
\end{egprob}
}{title=H18東工大}
\begin{egans}
指数$p$の部分群$H$を取る. cosetの集合$G/H$を考え, 置換表現$\rho:G\to \mr{Aut}(G/H)$; $gH \xmapsto{g'} g'gH $を考える. $\mr{Aut}(G/H)$の位数は$p!$である.
\begin{claim}
$H = \ker{\rho}$である. とくに, $H$は正規部分群である.
\end{claim}
[証明] $[G:H]=p$, $[G:\ker{\rho}] = |\im{\rho}|$である. \aka{ $k\in \ker{\rho}$であるとき, とくに$kH = H$であることから$k\in H$なので, $\ker{\rho} \subset H$が従う.}(これを見落としがちである)  よって$[G:\ker{\rho}] = [G:H][H:\ker{\rho}] = p[H:\ker{\rho}]$であるから$[H:\ker{\rho}]=1$を示せばよい. そこでこれが$[H:\ker{\rho}] > 1$であったとしよう. いま,  
\[\frac{1}{p}|\im{\rho}| = [H:\ker{\rho}]\]
が分かっている. $|\im{\rho}|$は$p!$の約数だから, $[H:\ker{\rho}]$は$(p-1)!$の約数である. よって仮定よりこの数は$p$未満の素因数$q$を持つ. $|G| = [G:H][H:\ker{\rho}][\ker{\rho}:\set{e}]$だから$|G|$の約数でもある. よって$|G|$は$q$で割り切れるがこれは$p$の最小性に反する. \qed 
\end{egans}



ガロア理論と絡めて出題されるようなこともある。正直, 嫌な問題だと思う。
\begin{egprob} (H1-I-3)\label{prob:H1_I_3} 
可換体$K$の5次拡大体$K(\theta)$で, $K(\theta)$を含む$K$の最小のGalois拡大体が$K$上の40次になるものは存在しないことを示せ。
\end{egprob}
\begin{screen}
やけに問題が一般的だと思わないだろうか？5次拡大体なんていろいろあるだろうに, このようなことが成り立つというのなら,  この現象が群論レベルの出来事である可能性が高い。
\end{screen}
\begin{ans}
$5$次対称群$\maf{S}_{5}$に指数$3$の部分群は存在するかを考えればよい。ここで使う非自明な事実は, $A_5$が単純群であるということである。\par 
\tf{指数3の部分群は存在しない。}これを背理法によって示す。

$H$を$\maf{S}_{5}$の指数$3$(位数$40$)の部分群とする。$40$は$60$を割りきらないから, $H$は$A_{5}$の部分群ではない。よって$H$には奇置換が存在し, 奇置換と偶置換がそれぞれ$20$個ずつ存在する($\because$奇置換$\sigma$を一つ選び, $H\ni g\mapsto g\sigma\in H$という$H$の自己同型が$H$の奇置換と偶置換の全単射を与えるから)。$K:=H\cap A_5$とすると, $K$は位数$20$である。cosetの集合$X=A_{5}/K$を考え, $A_5$の$X$への作用を$\sigma K \mapsto (\tau\cdot \sigma )K$で定める。つまり, 群準同型$\phi:A_5\to \maf{S}_{3}$を
\[\phi(\tau) = (f_{\tau}: \sigma K \mapsto (\tau \sigma )K)\]
で定義する。このとき$\ker{\phi}$は$A_5$の正規部分群であるが, $A_5$の単純性より$\ker{\phi} = \{ 0 \}, A_5$となるしかないが, 写像$A_5\to S_3$が単射になることはありえないので$\ker{\phi} \neq \{ 0 \}$であり, 一方$\ker{\phi} = A_5$だと$\phi$は自明な準同型だが $\phi$の定義から明らかに不適である。よって矛盾が得られた。\qed
\end{ans}



\subsection{演習問題}
\subsubsection{Level B}

$S_5$に指数3の部分群が存在しないことを以前用いたが,  $S_4$には存在するというのも有名である. 次はそれが題材になっている. 

\begin{prob} (H31東北大専門1)
$G=\maf{S}_{4}$の二つの元$\sigma = (1234)$, $\tau = (24)$で生成される部分群を$H$とする. 
\ben 
\item $H$が非可換であることを示せ. 
\item $G/H$の元のそれぞれに対して代表元を一つずつ求めよ. 
\item $K=\set{g\in G\mid \forall x\in g, (gx)H = xH}$は$G$の正規部分群であることを示せ. 
\item $G/K\cong \maf{S}_{3}$を示せ. 
\een 
\end{prob}

\subsubsection{Hint・Answer}
{\small 
\begin{probans}
(1): $\sigma\tau \sigma^{-1} = (\sigma(2)\sigma(4)) = (31)\neq \tau$より. \\
(2): $H$の位数が8であることを示す. まず, $G$の元に対する次の性質を考える. 
\lefba{ ($\ast$)
$g\in G$を全単射$g:\set{1,2,3,4}\to \set{1,2,3,4}$とみなすとき, $g(\set{2,4}) = \set{2,4}$であるかまたは$g(\set{2,4}) = \set{1,3}$である. 
}
$g_1,g_2,g\in G$が$(\ast)$を満たせば$g_1g_2$と$g^{-1}$もこれを満たすので, $(\ast)$を満たす元全体は部分群をなす. $\sigma,\tau$はこれを満たすので, $H$のすべての元はこれを満たす. また, $(\ast)$を満たす$g$は, 組み合わせ的に8通りしか存在しないので, $H$の位数は8の約数. $H$は$1,\sigma,\sigma^2,\sigma^3,\sigma\tau$の5つの元を含むので, $\sharp{H}=8$が分かる. $H$の元は$(\ast)$で特徴づけられるということが分かったので, $(12),(34)\notin H$である.
\lefba{ ここで $(12)H \neq (34)H$を示す. もし等号が成り立つなら$(12)\sigma(34) \in H$だが, これは$(23)$であり, $(\ast)$を満たさないので$H$に属さない. よって矛盾. 
}
よって$G/H$は3つの元を持ち, $1,(12),(34)$で代表される. \\
(3): $K=\bigcap_{x\in G}x^{-1}Gx$と書けるので$h^{-1}Kh = K$はすぐ分かる. \\
(4): $K$はcosetへの作用$\phi:G\to \mr{Aut}(G/H) \cong \maf{S}_{3}$の核である. よって$G/K\cong \im{\phi}$. 
\[\sigma^2 = (13)(24), \sigma\tau = (12)(34), \sigma^3\tau = (23)(14)\]
はサイクル分解で, 同じ型を持っているので$G$で共役である. (2)の代表元を用いると, 
\[(12)\sigma^2(12) = (14)(23) = (34)\sigma2(34) \in H\]
なので, $\sigma^2 \in K$である. $K$は正規だから$K$は$\sigma^2$の共役元も含み, 位数4以上であることが分かる. 一方, $(12)\tau(12) = (14)\notin H$なので$\tau\notin K$であり, $\sharp{K} = 4$が確定する. よって$G/K$も$\maf{S}_{3}$も位数が6なので$G/K\to \maf{S}_{3}$は全単射準同型であり, 同型である. \qed 
\end{probans}

}



\newpage 




\section{一般線形群・特殊線形群}

\subsection{有限体上の線型代数}
以降, $p$は素数とする. 有限体上で線型代数を考える問題が専門レベルではたまにある. 代数閉体でもないので, 何が有限体上で行えるのか, それとも代数閉体をつかわないと行けないのかに注意することは大事である. たとえば, $\mb{F}_{p}$成分の行列は必ずしも$\mb{F}_{p}$において固有値を持つとは限らないし, 対角化するときの行列$P$は$\mr{GL}_{n}(\ovl{\mb{F}_{p}})$の中に存在するが, $\mb{F}_{p}$成分であるとは限らない. \par 
次は覚えておきたい問題である. 
\begin{egprob}
$\mr{GL}_{n}(\mb{F}_{p})$の位数を求めよ. 
\end{egprob}
\begin{egans}
正則性と列ベクトルの一次独立性は有限体上でも同値である. したがって, $\mb{F}_{p}^{n}$の順序を考慮した基底の個数だけあるので, 組み合わせ論的に考えると
\[\bolm{(p^n - 1) (p^n - p) \cdots (p^n - p^{n-1})}\]
となる. 
\end{egans}


\begin{egprob}
$X\in \mr{M}_{2}(\mb{F}_{p})$であって, $X^p = I$なるものはいくつあるか. 
\end{egprob}
\begin{egans}
$\ovl{\mb{F}_{p}}$を一つ固定して$X^{p} = I$なる$X$をJordan標準形にする$P\in \mr{GL}_{2}(\ovl{\mb{F}_{p}})$を取る. $P^{-1}X^{p}P = I$なので, $P^{-1}XP = \begin{bmatrix}a&b\\ 0& d\end{bmatrix}$とすると$a^{p} = d^{p} = 1$である. $a,d\in \ovl{\mb{F}_{p}}$なので, $a^{p^r} = a$, $d^{p^s} = d$となる$r,s$がある. ここから$a=d=1$を得る. すなわち$X$の固有多項式は$(t-1)^2$である. 逆に$X$の固有多項式が$(t-1)^2$ならば, Jordan標準形において
\[P^{-1}X P = \begin{bmatrix} 1& b\\ 0&1 \end{bmatrix} \]
の形になるが, このとき
\[X^{p} = P \begin{bmatrix}1&b\\ 0&1\end{bmatrix}^{p} P^{-1} = \begin{bmatrix}1 & 0 \\ 0 & 1\end{bmatrix}\]
となる. よって, 固有多項式が$(t-1)^2 = t^2 - 2t + 1$であるような$X$をｓべて求めればよいから, 
\[a+d = -2,\quad ad-bc = 1\]
となるような$a,b,c,d$がいくつあるかを求めればよい. それは頑張ればできる. 
\end{egans}



\subsection{演習問題}

\subsubsection{Level A}

\subsubsection{Level B}

\subsubsection{Level C} 

\begin{prob} (H28専門科目) \label{prob:H28senmon-3}
$G = \mr{GL}_{2}(\mb{F}_{p})$とおく. 
\ben 
\item $G$の元で対角行列と共役でないものの個数を求めよ. 
\item $G$の二つの元$X,Y$の最小多項式$\phi_X(t)$, $\phi_Y(t)\in \mb{F}_{p}[t]$が一致すれば $X,Y$は互いに共役であることを示せ. 
\item 
\een 
\end{prob}



\subsubsection{Hint・Answer} 
\begin{probans}
(1): $\sumib{a&0\\0 &d}$と共役な対角行列は自身と$\sumib{d&0\\ 0&a}$のみ. $G$から$G$への共役作用を考え, 各$x=\sumib{a&0\\ 0&d}$の軌道$O_{x}$の元の個数を考え,次のように足せばよい. 
\[\sum_{a\in \mb{F}_{p}^{\times}} |O_{aI}| + \dfrac{1}{2}\sum_{a,d \in \mb{F}_{p}^{\times} } |O_{\sumib{a&0 \\ 0 & d}}|  \]
$|O_x||G_x| = |G| = (p^2-1)(p^2-p)$なので, $|G_x|$を求めると答えが分かる. さて, $P^{-1}(aI)P = aI$だから$|G_{aI}| = |G|$より$|O_{aI}| = 1$. また, $P=\sumib{x&y\\ z&w}$のとき, 
\[\bmat{a& 0\\ 0&d}P = P\bmat{a&0 \\ 0 &d}\implies \bmat{ax & ay \\ dz & dw} = \bmat{ax & dy \\ az & dw}\]
となるから, $a\neq d$ならば$y=z=0$が条件. よって$P$は$(p-1)^2$通りあり, $a\neq d$なら $|G_{\sumib{a&0\\0&d}}| = (p-1)^2$となる よって, $|O_{\sumib{a&0\\ 0&d}}| = p(p+1)$. 以上より求める個数は
\[\dfrac{1}{2}(p^2- 1)(p^2-p) - (p-1)\]
(2): 


\end{probans}








\clearpage  
\section{その他}
この接で述べることは代数学入門以降の内容である. 
\subsection{p群}
$p$群の顕著な性質は以下に代表される. 
\begin{prop}
$p$群$G$の中心は非自明である. 
\end{prop}
\prf 
類等式を考えて, 中心の位数が1だとすると, 中心に属さない元の共役類の位数がp冪で, 類等式の$\bmod{p}$の値で矛盾が起こる. 
\qed 

$p$群に関する一般的な問題に対しては中心の非自明性に必ず注意されたい. 具体的な$p$群の問題については, そもそも位数の小さい$p$群自体限られているものであるから, いっそのこと$p$群の例を覚えておくのがよい. 
\begin{eg} $p$は素数とする. 
\ben 
\item 位数$p$の群は巡回群$C_p$に限る. 
\item 位数$p^2$の群は必ずAbel群であり, 構造定理から$C_{p^2}$あるいは$C_p\times C_p$である. すべての元が位数$p$であると分かれば, $C_p$である. 
\item 位数$p^3$のAbel群は$C_{p^3}, C_{p^2}\times C_p, C_p^3$の3通り. 非アーベル群については次の2つがある. 
\lefba{($C_{p^2}\rtimes C_p$) 
$C_{p^2}^{\times}$には位数$p$の元$a$が存在する. そこで, $\phi: C_p\to \mr{Aut}(C_{p^2})$を$\phi(x) = (\phi_x:t\mapsto a^x t)$により定める. $a^{pm+x} = a^x$だから, これはwell-definedである. これにより$C_{p^2}$の$C_{p}$による半直積$C_{p^2} \rtimes_{\phi} C_{p}$が得られる. その演算は,
\[(n_1,h_1)(n_2,h_2) = (n_1+a^{h_1}n_2, h_1+h_2)\]
というものになる. 
}
\lefba{(\tf{Heisenberg群}; $\mr{Heis}(\mb{F}_{p})$)
$\phi:C_p\to \mr{Aut}(C_p\times C_p)$を$\phi(t) = (\phi_t:(a,b)\mapsto (a+bt, b))$により定める($\phi_t$が自己同型であることを確かめよ). これにより半直積$(C_p\times C_p)\rtimes_{\phi} C_p$を得る. この演算は, 
\bealn{ 
&((a_1,b_1),c_1)((a_2,b_2),c_2) \\
&= ((a_1,b_1)+\phi_{c_1}(a_2,b_2), c_1+c_2) \\
&= ((a_1,b_1) + (a_2+c_1b_2, b_2), c_1+c_2) \\
&= ((a_1+a_2+c_1b_2, b_1+b_2), c_1+c_2)
}
そして, この群はHeisenberg群$\mr{Heis}(\mb{F}_{p})$に同型である. この群は, $\mr{GL}_{3}(\mb{F}_{p})$の部分群として, よく次の形で記述される:
\[\mr{Heis}(\mb{F}_{p}) := \set{\bmat{1&c&a\\ 0&1&b\\ 0&0&1} \mid a,b,c\in \mb{F}_{p}}\subset \mr{GL}_{3}(\mb{F}_p)\]
この群では演算が
{\small
\[\bmat{ 
1&c_1&a_1 \\ 0&1&b_1\\ 0&0&1
}
\bmat{
1&c_2&a_2 \\ 0&1&b_2\\ 0&0&1
}
= \bmat{1&c_{1}+c_2 & a_1 + c_1b_2 + a_2 \\
0&1& b_1+b_2\\ 0&0&1}
\]
}
であり, 先に構成した$(C_p\times C_p)\rtimes_{\phi}C_{p}$と同型であることは視覚的に明らかである. \tf{特に, この群はすべての元が位数$p$である. }
}
\item ただし, 上二つで挙げた群は$p=2$のときに二面体群$D_4$に同型で, 位数$8$の非アーベル群はこのほかに$Q_8 = \set{\pm 1,\pm i, \pm j ,\pm k}$が存在する. 
\een 
\end{eg}

\begin{prac} 
(3)(a)の例において, $a$の取り方には依存するが, $a'$を別の位数$p$の元として取り替えて半直積を構成した場合, それらは同型になることを示せ. (未検証)
\end{prac}


\begin{egprob} (H21数学II-2) 
$p$を奇素数とする. $G$が位数$p^3$の有限群で, 単位元以外の各元の位数$p$であるようなものとする. このとき, $G$は$\mr{GL}_{2}(\mb{C})$の部分群と同型ではないことを示せ. 
\end{egprob}
\begin{egans}
部分群$H$と同型とせよ. $X\in \mr{GL}_{2}(\mb{C})$で$p$乗して単位行列になるものについて考えよう. 明らかに$X$の固有値は$\zeta_p^{k}$と表せる. もし$X$が対角化可能なら, $X$は対角成分の元が$\zeta_{p}^k$であるものと相似なので, $X^{p} = I$が従う. 逆に$X^p = I$であるなら, $T^p - 1$は重根を持たないので, 最小多項式もそうである. よって$X$が対角化可能になる. 
$P\in \mr{GL}_{2}(\mb{C})$のとき,  $H\cong P^{-1}HP$なので$H$は単位行列以外の対角行列を含むとしてよい. それを$A=\mr{diag}(\zeta_p^i, \zeta_p^j)$とする. ここで$\zeta_p^j \neq 1$なら, $kj\equiv 1\pmod{p}$なる$k$を取って$A^k$を考えて取り変えることで$\zeta_p^j = \zeta_p$としてもよい. もし$\zeta_p^j = 1$なら, $\zeta_p^i \neq 1$だから, $S = \sumib{0&1\\ -1&0}$として$H\cong S^{-1}HS$に注意することで$\zeta_p^j \neq 1$の場合に帰着する. 結局, $A=\mr{diag}(\zeta_p^i, \zeta_p)\in H$であると仮定してよい. 

$Z(H)$は非自明なので, この中心を解析する. $B=\sumib{a&b\\c&d}\in Z(H)$とせよ. 
$AB = \sumib{\zeta_p^i a& \zeta_p^i b\\ \zeta_p c & \zeta_p d}$, $BA = \sumib{\zeta_p^i a & \zeta_p b \\ \zeta_p^i c & \zeta_p d}$だから$i=1$であるか, $i\neq 1$で$b=c=0$である. \par 
いずれの場合も, $Z(H)$は対角行列から成る. 単位元でないものが取れるので,そのひとつを$B = \mr{diag}(\zeta_p^{k}, \zeta_p^{l})$とする. \par 
位数$p$の対角行列は$p^2$個しかないので, 対角行列ではない$C\in H$が必ず存在する.  $C=\sumib{s&t\\ u&v}\in H$ で$CB = BC$を解くと\ao{$k=l$でなければならないことが分かる.} よって$B_1=\mr{diag}(\zeta_p,\zeta_p) \in Z(H)$となることも分かる. $B$は$Z(H)$の単位元でない任意の元であったから, $Z(H)$は$B_1$の冪からなり, 特に$\sharp Z(H) = p$が分かる. $H$を$P^{-1}HP$に取り替えるにあたって, この中心は変わらないことに注意すると, もともと$B_1$の冪が$H$に含まれていることが分かる. 先の$C$を対角化する行列$Q$を取り$Q^{-1}CQ\in Q^{-1}HQ$を考える. $Q^{-1}CQ$はこのとき$B_1$の冪にはならない対角行列である. よって$H$に$C_0=\mr{diag}(\zeta_p^{k}, \zeta_{p}^{l})\notin Z(H)$が含まれるとしてよい. $B_0$と$C_0$をいくつか掛けていくことで, 任意の位数$p$の対角行列を作ることができる. $K$をこの対角行列全体のなす部分群とする. $C = \sumib{a&b\\ c&d}$が対角行列でないとすれば, $T_{k}=\mr{diag}(\zeta_p^{i_k}, \zeta_p^{j_k})\in K$に対し$T_{1}CT_{2} = \sumib{\zeta_{p}^{i_1+i_2}a & \zeta_{p}^{i_1+j_2}b \\ \zeta_{p}^{i_2+j_1}c & \zeta_{p}^{j_1+j_2}d}$である. $C$は対角行列でないが, $a=d=0$では$C^p \neq I$なので$C$の成分で0のものは1個以下である. このとき, $i_1,i_2,j_1,j_2$を$0$から$p-1$以下で動かすとき, $T_1CT_2$が$p^3$個以上の異なる(対角でない)行列を作ることが分かるが, 単位行列と合わせると$H$の位数を越えるので, 矛盾である. \qed 
\end{egans}

\begin{egans} (他の方から教えて貰った有限群の表現論による解法)
埋め込み$\rho:G\to \mr{GL}_{2}(\mb{C})$を2次元表現と見る. 次を認める. 
\begin{prop}
一般の群$G$について, 既約$G$表現の次元は$|G|$の約数である. 
\end{prop}
これによると$2$は$p^3$の約数にならないので$\rho$は既約表現になり得ない. よって可約であり, $\rho$は自明表現二つの直和に同型で, $\rho$は$\mb{C}^{\times}\times \mb{C}^{\times}$の部分群に同型である. それは$\cyc{p}\times \cyc{p}$の部分群に同型であり, 矛盾. \qed 

\end{egans}

\begin{egprob}
(位数$p^3$の群の分類を用いずに)位数$p^3$の群$G$は$|Z(G)| \neq p$を満たすことを示せ.
\end{egprob}
\begin{egans}
$|Z(G)=p^2$であるとしよう. $G/Z(G)$は巡回群であり, これの生成元を$b$とする. このとき任意の$G$の元は$b^kz$ ($z\in Z(G)$)で表せる. このような元を二つ用意し, $b^{k_1}z_1, b^{k_2}z_2$とする. 
\[b^{k_1}z_1b^{k_2}z_2 = b^{k_1+k_2}z_1z_2\]
だから$G$はアーベル群である. $|Z(G)| = p^3$なので, 矛盾. よって$|Z(G)|=p^2$ではない. 
\end{egans}




\subsection{可解群}
R2の問題では可解群が登場した. 
\begin{defn}
$G$が可解群であるとは, 部分群の列$\set{1} = H_0 \subsetneq H_1 \subsetneq \dots \subsetneq H_{n-1} \subsetneq H_n = G$が存在して, 
\ben 
\item $H_i \norgr H_{i+1}$ ($0\leq i < n-1$)
\item $H_{i+1}/H_i$はAbel群である. 
\een 
\end{defn}

\begin{rem}
定義から$H_1$は部分Abel群でなければならない. 
\end{rem}

正規部分群で繋げなければならないので, 中心$Z(G)$や指数2の部分群であるなど, 一般的に正規部分群と分かる群のことを抑えておくのがよいだろう. 


\begin{theo}
$p$群は可解群である. 
\end{theo}
\prf $G$の位数$p^n$として$n$に関する帰納法で調べる. $n=1$の場合は明らか. $n-1$までの主張を仮定する.  $Z(G)$は非自明で$G$の正規部分群だから, $G/Z(G)$は$p$群で位数$p^{n-1}$以下. よって$G/Z(G)$は帰納法の仮定より可解. $Z(G)$は可換だから可解. 
\qed 


\tbox{ 
\begin{point}
\ben 
\item \tf{Sylow $p$ 部分群$P$を考察せよ.} もしそれが正規($n_p = 1$)なら, まず$P\norgr G$. もしくは$G/P$が正規と言えればよい. 
\item $N$は正規部分群で, $N$と$G/N$がともに可解群であるようなものが存在するとする. このとき$G$も可解群である. 
\een
\end{point}
}{title=テクニック}


\begin{egprob}
位数1998の群$G$は可解群であることを示せ. 
\end{egprob}
\begin{egans}
$1998=2 \times 3^3 \times 37$. Sylowの定理よりSylow 37 部分群$S_{37}$が一つあり, それは正規. よって$G/S_{37}$が可解群と言えばよい. $G' := G/S_{37}$のSylow $3$部分群$S_{3}$は, 指数2なので正規部分群. よって$G'/S_3$は明らかに可解群. よって$G$は可解群.
\end{egans}

\begin{prac}
位数$2p^nq^m$ ($p,q$は素数, $2p^n < q$) の群は可解群であることを示せ. 
\end{prac}

\begin{egprob}
位数275の群は可解群であることを示せ. 
\end{egprob}
\begin{egans}
$275 = 5^2 \times 7 \times 13$. $S_{13}$は可換かつ正規なので$G' = G/S_3$の可解と言えばよい. $G'$のSylow 5部分群$S_{5}$は正規かつ可換(位数$p^2$の群は可換)で $G'/S_{5}$も位数7で可換なので$G'$は正規. よって$G$は正規.   
\end{egans}

\begin{egprob}
位数2020の群$G$は可解群であることを示せ. 
\end{egprob}
\begin{egans}
$2020 = 4\times 5\times 101$である. $S_{101}$は正規で, 可換より可解なので$G' = G/S_{101}$が可解であればよい. $G'$の$S_{5}$は正規で, 可換より可解なので$G'/S_{5}$が可解であればよい. これは位数$2^2$なので可換であり, 可解である. 
\end{egans}

\begin{egprob}
$p,q$が相異なる素数であるとき, 位数$pq$の群$G$は可解であることを示せ. 
\end{egprob}
\begin{egans}
$p<q$としてよい. $S_q$は正規で, 可換なので可解である. よって$G/S_{q}$が可解と言えばよいが, これは位数$p$なので可換とくに可解である. よって$G$も可解である. \qed 
\end{egans}

次は一筋縄ではいかない場合である. 院試レベルとしては次が出来るようになるのがよい. 
\begin{egprob}
$p$が素数であるとき, 位数$4p$の群$G$は可解群であることを示せ. 
\end{egprob}
\begin{miss}
$p\geq 5$のとき, Sylow $p$部分群$S_p$は正規でAbel群なので$G':=G/S_p$が可解であればよい. $G'$は位数4だから, これもAbel群でありよい. \par 
$p=3$のとき, $S_3$が正規なら先と同様にすればよい. 正規でないこともある(\tf{ここで諦めてはならない！！！}).  正規でないとせよ. このとき$n_3 = 4$であり, 位数3の元が$8$個以上ある. このときSylow 2部分群$S_2$が正規にならねばならない. [証明] もし$n_2 > 1$なら$n_2 = 3$であり, \aka{位数2または4の元が$9$個以上ある}ことになり, $8+9  > 12$だから矛盾. [証明終] よって$S_2$は正規かつ(位数$p^2$だから)可換で可解. そして, $G/S_2$は位数3だから可解. よって$G$も可解. \par 
$p=2$のとき, $G$はp群であるから可解である. あるいは, 次のように直接調べる: $Z(G)$が非自明である. $Z(G) = G$なら$G$は可換なのでよい. $Z(G)$の位数2,4なら$Z(G)$は可換で可解かつ$G/Z(G)$の位数が4,2なので可換で可解. よっていずれの場合も可解. 
\end{miss}
\begin{egans}
赤線の部分は誤りである(一方で位数3の元が8個あることは正しい！). 9個以上とまではいえない. なぜなら,$S_{2,1},S_{2,2}$がSylow-2部分群, つまり位数4のとき, $S_{2,1}\cap S_{2,2} \neq \set{e}$となるかもしれないからだ. \par 
そこで次のように考えるとよい. $S_{2,i}$ ($i=1,2,3$)がSylow-2部分群とする. このとき様々な組み合わせを考えると$S_{2,1}\cup S_{2,2}\cup S_{2,3}$が7個以上の元を持つことが分かり, 位数1,2,4の元が7個以上あることが分かり, 矛盾である. \qed 
\end{egans}

\begin{egprob}
位数$2p^n$の群は可解であることを示せ. 
\end{egprob}
\begin{egans}
$S_p$は指数2なので正規で可解よりすぐわかる. 
\end{egans}

\begin{egprob}
位数$2pq$ ($3\leq p<q$)の群は可解であることを示せ. 
\end{egprob}
\begin{egans}
$n_q = 1$の場合は$S_q$が正規で, 位数$2p$の群は可解なのでよい. $n_q\equiv 1\pmod{q}$かつ$n_q \in \set{1,2,p,2p}$であるが, $n_q \neq 1$だとすると$n_q \geq q-1 >  p$なので$n_q = 2p$である. この場合$2p-1 = qk$と書けるが$p<q$により$k=1$である. よって$2p=q+1$となる. \par 
$n_q = 2p$のとき位数$q$の元が$2p(q-1)$個だけあり, そうでないものが$2p$個あると言える(この中に単位元もある). 
一方$n_p \in \set{1,2,q,2q}$である. $n_p = 1$のときは$G/S_p$が可解なのでよい. $n_p > 1$のときを考える. $n_p(p-1)$個の元が位数$p$であることが分かるから, $n_p(p-1) \leq 2p-1$でなければならない. よって$n_p = 2$である. このとき位数$p$の元は$2p-2$個はあるといえる. よって, 位数$p,q$でない元が単位元を除きあと1個ある. \par 
一方$n_2 \in \set{1,p,q,pq}$であり, $n_2 > 1$だと位数2の元が$2$個以上あることになるので現在の状況に矛盾する. よって$n_2=1$で$S_2$が正規. $G/S_2$は可解なのでよい. \qed 
\end{egans}
\begin{egprob}
位数$56$の群が可解であることを示せ. 
\end{egprob}
\begin{egans}
$n_7$の可能性は1と8. $n_7 = 1$ならよい. $n_7 = 8$なら48個以上, 位数7の元が存在することが保証される. もしこのとき$n_2 = 7$なら位数2の元が$8$個以上(実際はもっと)あることが保証される. これと単位元をあわせると$56$を越えるのでおかしい. よって$n_2 = 1$なので$G/S_2$が可解と言えばよく. これは位数7なのでOK. \qed 
\end{egans}




\begin{egprob}
位数$p^2q$ ($p<q$)の群は可解であることを示せ. 
\end{egprob}
\begin{egans}
$n_q=1$のときは$G/S_q$が位数$p^2$で可換. $n_q>1$とする. $p-1$は$q$で割れないので, $p^2-1 = (p-1)(p+1)$が$q$で割り切れ, $p+1$が$q$で割り切れなければならない. $p < q \leq p+1$となる. よって$q=p+1$だが$p,q$が素数なので$p=2,q=3$となる. この場合は, $4p$のときの考察で完了している. \qed 
\end{egans}

次も知識として知っておくとよい. 
\begin{theo}
位数59以下の群は可解群である. 
\end{theo}
\prf まず, $1,\dots,59$を次のパターンに分類する. 
\ben 
\item 素数の累乗: 1, 2, 3, $2^2$, 5, 7, $2^3$, $3^2$, 11, 13, $2^4$, 17, 19, 23, $5^2$, $3^3$ ,29,31, $2^5$ ,37, 41, 43, 47,$7^2$, 53, 59. このときは$p$群なのでよい. 
\item 異なる二つの素数の積: $6,10,14,15,21,22,26,33,35,38,39,46$, $51,55,57,59$. 
\item $4p$: $12,20,28,36,44,52,$
\item $2p^n$: $18,50,54$
\item $2pq$: $30,42$
\item $p^2q (p<q)$: 45がこれに当てはまるのでよい. 
\item それ以外: $24,40,48,56$
\een
24,56についてはすでに調べた. 
このうち, 位数が40のものが可解であることは簡単にわかる. 
\begin{prac}
位数40の群が可解であることを示せ. 
\end{prac} 
%8 times 5で n_5 = 1なのでOK 
少し難易度が上がるが24も簡単である. 
\begin{prac}
位数24の群が可解であることを示せ. 
\end{prac}
%\begin{egans} $n_2=1$ならよい. $n_2 = 3$なら位数2冪の元が$21$個はあり, このとき$n_3 \geq 3$にはなりえない. よって$n_3=1$で$G/S_3$が可解なのでよい. 
%\end{egans}
よって次を示せばよいが, これは難易度が高い. 
\begin{claim}
位数48の群は可解である. 
\end{claim}
\prf $S_2$が正規の場合はよい. $S_2$が正規でないとき, $n_2 = 3$であり, それらを$S_{2,i}$ ($i=1,2,3$)とする. ここで\tf{シロー部分群は互いに共役であるから}$N:=S_{2,1}\cap S_{2,2} \cap S_{2,3}$は正規部分群である. もし$N$の位数が2以上なら, $G/N$が位数24,12,6,3のどれかなので$G$も可解であることが分かる. $N=\set{1}$とする. この場合, 実は異なる$i,j$について$S_{2,i}\cap S_{2,j} = \set{1}$であることを示そう. $i=1,j=2$としてよい. $x\in S_{2,1}\cap S_{2,2}$とする. このとき\bolm{x^{-1}S_{2,3}x = S_{2,k}}となる$k$があるが, これは$k=3$でなければならない. よって$x\in N_{G}(S_{2,3})$である. ところで, Sylowの定理より$N_{G}(S_{2,3})$の指数が$n_2 = 3$で, 正規化群の定義より$N_{G}(S_{2,3})\supset S_{2,3}$なので, もし$x\neq 1$なら$x\notin S_{2,3}$だが, このとき$N_{G}(S_{2,3})\supsetneq S_{2,3}$となるのに,  $S_{2,3}$と$N_{G}(S_{2,3})$がともに指数3であるのはおかしい. よって矛盾で, $x=1$だから$S_{2,i}\cap S_{2,j} = \set{1}$である. よって, $G$には位数2,4,8,16の元が少なくとも$45$個ある. もしこのとき$n_3 \geq 2$なら, 位数3の元が少なくとも4個あり, おかしい. よって$n_3 = 1$であり, このとき$G/S_3$が可解だからよい. \qed 
\begin{prac} 
上の証明をガロア理論から解釈せよ. 
\end{prac}



次の定理は知識用である. 答案には使えないであろうが, 覚えておくといいかもしれない. 
\begin{theo} (\tf{Burnside}の定理) 
$p,q$を素数とする. 位数$p^m q^n$の群は可解である. 
\end{theo}
よって, 位数$2021 = 43\times 47$の群は可解群である. また, 位数$4p$の群はやはり可解である. 
\begin{prac}
位数2021の群は可解であることをBurnsideの定理を使わずに示せ. 
\end{prac}

最後に, 位数60以降の話を少しだけ見てみよう. 
\begin{prop}
位数$\dfrac{n!}{2}$ ($n\geq 5$)の群は可解とは限らない. 
\end{prop}
\prf $A_n$($n\geq 5$)は単純群だから. \qed 

\begin{egprob}
$S_{n}$ ($n\geq 5$)が可解群ではないことを示せ. $A_n$が単純群であることは用いてよい.  
\end{egprob}
\begin{egans} $S_n$の正規部分群$N\neq \set{1}$を考える. このとき$N\cap A_n\norgr A_n$より$N\cap A_5 = \set{1}$でなければならない. どのような部分群も偶置換は必ず含むので, $N=A_n$となるしかない. よって系列は$A_n\norgr S_n$で終わらなければならず, そのあと$A_n$の部分群の列を作らなければならないが, これは$A_n$の単純性から不可能. \qed 
\end{egans}


\begin{prac}
位数61,62,63,64,65,66,67,68,69,70の群は可解であることを示せ.
\end{prac}





\subsection{演習問題}





\clearpage 





\section{京大過去問集}
\subsection{基礎問題}
\begin{prob} \checkbox (H28基礎科目II) \\
$A = \bmat{0&-1\\ 1&1}$, $B = \bmat{0&1\\ 1&0}$で生成される$GL_2(\mb{C})$の部分群$G$について以下の問に答えよ. 
\ben 
\item $|G|$を求めよ. 
\item $|Z(G)|$を求めよ. 
\item $G$の位数2の元の個数を求めよ. 
\een
\end{prob}
\begin{point}
生成元が二つあたえられているので, とりあえずいろいろ計算して関係式を出してみる. 共役元を求めるといいことがたまにある.  (1)の問題の出し方的に, $A,B$は有限位数であることにも注意. 
\end{point}
\begin{ans}
$A^6 = B^2 = I$はすぐ分かる. $BAB^{-1} = \sumib{1&1\\ 0 & -1}\sumib{0&1\\ 1&0} = \sumib{1&1\\ -1&0} = A^{-1} = A^5$である. よって$BA = A^{-1}B$. $G$の元は, $A^6 = B^2 = I$に注意すると
\[A^{k_1} BA^{k_2}B\dots A^{k_l}B^{b},\quad (0\leq k_i\leq 5, b=0,1)\]
の形に書けるが, $BA= A^{-1}B$を用いると必ず$A^{a}B^b$の形に書ける(ちゃんとやりたいなら$l$に関する帰納法を用いる). ただし$0\leq a\leq 5, b=0,1$. よって, この表記法が12通りしかないので$|G| \leq 12$がわかる. 一方で$A$の位数が6なので$|G|$は6の倍数であり, $B$は$A$の冪にならないので$|G|$は6より大きい. よって$|G| = 12$が従う. \\
(2): $|G|=12$より, $G$の元を$A^aB^b$ ($a=0,\dots, 5, b=0,1$)の形に表す方法は一意. 実は, このような表記に関して$b=1$であるものは中心に属さない. つまり, $A^aB\in Z(G)$となることはない.  なぜならば, もし$A^aB\in Z$なら
\[A(A^a B) = A^{a+1}B = (A^a B)A = A^{a-1}B\]
となり表記が一意であることに反するから. 次に $A^a\in Z(G)$であるための必要条件を考えよう. $BA=A^{-1}B$より$ABA = B$なので, $B=ABA = A(ABA)A = A^2BA^2 = \dots = A^nB A^n$である. よって, 
\[A^a B = BA^a = A^{-a} (A^a B A^a) = A^{-a} B \]
となり, 表記の一意性から$a\equiv -a\pmod{6}$である. よって$a=0,3$が必要. 逆にこのとき$A^a = \pm I$なので明らかに中心に属す. よって$|Z(G)| = 2$. \\
(3): $A^aB$の形の元は位数2である. 実際, $A^a B \neq I$で
\[(A^aB)^2 = (A^aB A^a)B = B B = I\]
1であるから. よって$A^aB$ ($a=0,1,2,3,4,5$)は位数2. また, $A$は位数6なので, $A^a$の形の元で位数2のものは$A^3$のみである. 以上より$7$個ある. 
\end{ans}



\subsection{専門問題}
\begin{prob} \checkbox (R2専門科目)
位数$3p^n$ ($p$は素数) の群は可解であることを示せ. 
\end{prob}

\begin{prob} \checkbox (H29専門科目)
$G$を有限群とする. 
\ben 
\item $|\mr{Aut}(G)| = 1$のとき, $G$は自明群であるか, または$\mb{Z}/2\mb{Z}$と同型であることを示せ. 
\item $|G|$が奇数で$|\mr{Aut}(G)| = 2$であるような$G$をすべて求めよ.
\een 
\end{prob}

\begin{prob} \checkbox (H28専門科目)\\
$G = GL_2(\mb{F}_{p})$とおく. 
\ben 
\item $G$の元で対角行列と共役でないものの個数を求めよ. 
\item $X,Y\in G$の最小多項式が一致すれば$X,Y$は互いに共役であることを示せ.
\item 対角行列を含まない$G$の共役類の個数を求めよ.
\een 
\end{prob}

\begin{prob} \checkbox (H27専門科目)
$G$は非可換群で「$N_1,N_2\subsetneq G$が相異なる$\set{1},G$以外の正規部分群なら, $N_1\not\subset N_2$」という条件を満たすとする. 
\ben 
\item$N_1,N_2$が相異なる$\set{1},G$以外の正規部分群なら, $G=N_1\times N_2$であることを証明せよ. 
\item $G$の自明でない正規部分群の数は高々2個であることを証明せよ. 
\een 
\end{prob}

\begin{prob} (H26専門科目)
$K\subset \mb{C}$を部分体, $p$を素数とする. $\mb{C}$に含まれる任意の有限次拡大$L/K$に対し, $L=K$でなければ$[L:K]$は$p$で割り切れると仮定する. このとき, $\mb{C}$に含まれる任意の有限次拡大$L/K$に対し$[L:K]$は$p$の冪(1を含む)であることを証明せよ. 
\end{prob}








